\section{改进算法与独立评价协议}
\subsection{快速组装与自适应回退}
固定网格与热源后，几何项 $K,D_s,b$ 可以预存。完整模型的公平快速对照每次只更新四个边界系数，然后仍完整求解同一稀疏线性系统。降阶模型则预存
\[
 K_r=V^\top KV,\qquad D_{r,s}=V^\top D_sV,\qquad b_r=V^\top b,
\]
在线求解 $(K_r+\sum_s\beta_sD_{r,s})a=b_r$，再重构完整温度场 $Va$。参数分离属于标准投影降阶技术，本研究的贡献是可核查的实现、失效分析与成本比较，不宣称这一公式的新颖性。

自适应模型先作快速降阶预测 $u_r$，再计算积分方程残差 $b-Au_r$ 及解析屏障指标 $\eta$。若 $\eta\leq\tau$，返回预测场；否则\textbf{实际执行}快速完整求解并返回其结果。回退解随后作两遍整体正交投影消去，即两遍重正交（CGS2）；当其剩余二范数与完整解二范数之比大于 $10^{-10}$、当前秩小于 $\min(32,n^2)$ 时，加入归一化剩余向量并更新全部低维缓存。达到秩上限或向量相关时仍执行所需回退，但不增广。零解不增广，域外参数及非有限数值被拒绝。

原生开发程序执行了 653 项检查：原与快速完整模型、原与快速投影的一致性，残差和正交性，触发规则，以及零向量、相关向量、非法参数和最小网格等分支均通过。52 个开发参数上的两种完整解之差为零，两种固定降阶场的最大差为 $7.49\times10^{-16}$。按 12 个调参案例再接 16 个边界角点的顺序运行时，发生 15 次真实回退及增广，秩由 8 增至 23；这 15 次均属于回退前近似已达标的保守触发。该结果仅作开发证据，不能当作独立测试表现。

\subsection{测试隔离与五种对照}
评价协议在最终测试生成前固定：每边网格数 $n\in\{24,48,96\}$，每个网格沿用相同的 24 个训练参数、重新计算对应网格上的训练解与奇异值分解。初始秩统一使用开发阶段选定的 $r=8$，不再根据最终测试调秩。测试包含 256 个按既定顺序使用的参数，随机种子为 20261027，各分量在 $[0.2,5]$ 内独立对数均匀采样。采样前保存算法、驱动与配置的摘要；模型函数不能取得用于事后评价的参照解。

比较五种方法：原始完整组装 FOM、预存几何的快速 FOM、完整组装再投影的固定 POD、预存低维矩阵的快速固定 POD，以及带屏障触发和基更新的自适应 POD。\textbf{主要速度对照使用快速 FOM}，以免把消除 Python 组装开销误算成降阶本身的收益。各方法均返回完整温度场，固定 POD 与自适应 POD 使用同一初始基，以区分组装优化和误差控制的作用。

场误差定义为 $\|u-u_{\mathrm{FOM}}\|_\infty$，热点值误差定义为 $|\max u-\max u_{\mathrm{FOM}}|$，并另记录相对二范数误差。误触发严格定义为 $\eta>\tau$ 且\emph{回退前}场误差不超过 $\tau$；不能用回退后接近零的误差判断。普通双精度计算的指标不具备区间运算的向外舍入保证，报告将精确算术误差界与实际浮点观测分开陈述。

\subsection{完整成本与来源归属}
对查询数量 $Q\in\{24,256\}$，定义每种独立部署方案的总成本
\[
 C_m(Q)=C_{m,\mathrm{offline}}+\sum_{i=1}^{Q}t_{m,i},
 \qquad S_m(Q)=\frac{C_{\mathrm{fast\ FOM}}(Q)}{C_m(Q)}.
\]
只有 $S_m>1$ 才代表该口径下的总成本降低。离线成本按方法的实际需要归属：原始 FOM 不承担训练或 SVD；快速 FOM 承担所需几何和源项预处理；各 POD 方案承担 24 次训练求解、SVD 与初始基准备，快速及自适应方案另计所需投影缓存，自适应方案还计屏障准备。每个网格、每次计时重复中可以只实际执行一次共享训练与 SVD，再复制同一初始基供对照，但在每种独立部署的 ROM 成本账上均完整计入一次。\emph{实际实验调用总数不因此乘以方案数量}。短批量也承担全部必要离线成本。自适应在线服务时间包含残差与指标、触发后的真实完整求解、重正交及缓存更新。普通 FOM 和固定 POD 不强制承担可选残差诊断，其诊断成本单列；事后误差评价不计入服务时间。

评价预定每个网格做三次计时重复，逐次重建训练解并重置基、缓存、状态与计数；实际完成次数以结果文件为准。重复使用同一冻结参数序列，是时间波动观测，不是新增独立精度样本。程序在导入数值库前将相关线程环境变量设为 1，但这本身不是对底层实际线程数的独立测量；环境设置与计时范围保存在产物中。数值服务时间也不包含大模型思考、工具调度和人工复核时间；后者属于研究开发成本，不应隐去或混作数值加速。

研究求解器、快速组装内核及自适应策略实现由原版 Moonshine 的工具调用生成并迭代；最终采样、计时、记账与结果汇总驱动由外部协调者提供，另行记录来源及摘要，再交 Moonshine 审阅与执行。独立数学检查、图表和报告整理同样属于外部复核与呈现。文件被复制到智能体工作目录，不改变其作者归属。
